\chapter{Las tres tesis}

\label{cap:planteo}
\section{Primera tesis: dos circuitos, regulación inversa}

En el departamento La Caldera operan dos circuitos de producción de suelo.

El primero es privado. Una parcela madre de gran superficie, generalmente de origen sucesorio o institucional, se estructura mediante fideicomiso o convenio, se desmonta, se mensura, se aprueba, se comercializa. Produce lotes de seiscientos metros cuadrados o más, para compradores de sectores medios y altos, en el lapso de un ciclo comercial.

El segundo es estatal. La Provincia adjudica tierra fiscal a familias de bajos ingresos, por sorteo público o por regularización de ocupaciones consolidadas, con un régimen de inscripción, puntaje, tenencia precaria y adjudicación en venta financiada. Produce lotes de interés social en plazos que, en los casos documentados, se miden en años y, en algún caso, en décadas: en un barrio de la capital en el que, según su centro vecinal, se adjudicaron lotes a vecinos del departamento, hay un expediente iniciado en 1995 y resuelto en 2025.

Los dos circuitos están regulados en proporción inversa a la capacidad de sus destinatarios. Al solicitante de un lote fiscal se le exige certificado policial de residencia, declaración jurada bajo pena de falso testimonio, referencias vecinales con nombre y apellido, franja horaria de visita domiciliaria, informe socioambiental y construcción en tres meses bajo apercibimiento de revocación. Al desarrollador que convierte monte de yungas en doscientas cuarenta y siete parcelas sobre una microcuenca judicialmente condenada se le exigen, en el papel, diez actos municipales ---los que enumera el Código de Ordenamiento (capítulo \ref{cap:cot})--- y, entre ellos y los provinciales, dos que aquí importan: un certificado de aptitud ambiental municipal y un \textbf{Certificado de No Inundabilidad} que la reglamentación del Código de Aguas califica de \textbf{recaudo indispensable} para todo loteo cuyo límite sea un río. \textbf{Ninguno de los dos aparece nombrado en el acto que aprobó el loteo más discutido del departamento}, y en el resto de los expedientes relevados sólo consta el de no inundabilidad, una vez, junto a un certificado ambiental provincial ---no el municipal---, en la respuesta de un loteador a una intimación municipal (capítulo \ref{cap:amparo}).

No se trata de un régimen laxo que favorece la discrecionalidad. Se trata de un régimen exigente aplicado a quien menos tiene y de un régimen tenue aplicado a quien más puede, sobre suelo del mismo corredor y bajo la misma autoridad provincial.

\section{Segunda tesis: la opacidad está en el criterio, no en el acto}

La segunda tesis no formaba parte del proyecto. Se impuso al cabo del relevamiento, cuando el mismo obstáculo apareció cinco veces seguidas en instrumentos de dos niveles de gobierno distintos, el municipal y el provincial.

La Ordenanza 437 de Vaqueros aprueba el Reglamento de Fraccionamiento de Tierras; el reglamento va en un anexo que el Boletín Oficial no publicó. La Resolución 82/16 aprueba el Reglamento de Asignación de Lotes de Interés Social y a él se remite el plazo de exhibición de los padrones; el reglamento va en un anexo que el Boletín sólo ofrece como imagen escaneada, que ningún buscador encuentra. La Resolución 8/16 establece el sistema de puntaje por vulnerabilidad social; el baremo va en un anexo que también se publica sólo como imagen. La Resolución 9/16 regula el sorteo público y su publicidad; el procedimiento va en otro anexo publicado del mismo modo. El Decreto 63/25 adjudica ciento tres parcelas conforme a un anexo que se publica aparte, como archivo vinculado, y que después se corrige por una resolución que subsana errores materiales.

El acto se publica. El criterio no se publica, o se publica de un modo que nadie lo encuentra: cuatro de esos cinco anexos resultaron estar en el repositorio del Boletín ---tres de ellos como imágenes escaneadas, sin texto que un buscador pueda leer---, y el quinto, el reglamento de Vaqueros, que el Boletín nunca publicó, apareció al final en el digesto en línea del propio municipio (capítulo \ref{cap:vaqueros}).

A eso se suma un mecanismo distinto y complementario. El Código de Ordenamiento Territorial de La Caldera declara sectores no urbanizables y, a continuación, prevé la figura del Proyecto Especial que permite urbanizarlos con más pasos. El proyecto de ordenanza del Plan de Desarrollo Urbano Ambiental emplea la fórmula ``quedará a consideración del municipio'' seis veces sin definir qué debe considerarse, con qué dictamen ni en qué plazo. La calificación de no urbanizable no opera como prohibición: opera como requisito de procedimiento agravado.

Ninguna de estas prácticas es ilegal. Los anexos existen, integran los actos y son accesibles por pedido; la excepción está prevista en la norma y se tramita conforme a ella. Y precisamente por eso el resultado es más eficaz que cualquier ocultamiento: no deja nada que denunciar.

El capítulo \ref{cap:opacidad} desarrolla esta tesis con el inventario completo. Se anuncia aquí porque el lector la va a encontrar antes, dispersa, en cada capítulo donde el libro se acerca a una decisión y encuentra publicada la envoltura y reservado el contenido.

\section{Tercera tesis: el dispositivo es anterior al mercado}

La tercera apareció como la segunda: por acumulación, y ya avanzado el trabajo. Surgió al
retroceder el relevamiento hasta el primer número del Boletín Oficial, en 1908.

Las unidades territoriales sobre las que hoy se lotea son las divisiones administrativas
del departamento a comienzos del siglo XX. Varios de los trece partidos policiales de 1909 —San Alejo, Los Porongos, Potrero de Gallinato, Potrero de Castillo, La Calderilla, Los Yacones— llevan los nombres de fincas y parajes que el proyecto de ordenanza del plan urbano de 2015 enumera como integrantes del municipio ---este libro leyó el proyecto, no el texto sancionado (capítulo \ref{cap:pdua})---; otros —Lesser, Vaqueros— nombran hoy parajes, cursos de agua o un municipio entero; y una fracción de Los Porongos es una de las que una ley de 1972 declaró afectadas. La finca Potrero de
Gallinato es partido en 1909, se remata judicialmente en 1922, la Provincia adquiere una
fracción por prescripción en 2015 y ese mismo año se aprueba sobre ella un club de campo de
mil quinientas sesenta y nueve hectáreas. La finca Cerro de Buena Vista se deslinda ante un
juez en 1921, obtiene agua en 1924 y llega a audiencia pública ambiental como urbanización
en 2014.

Lo mismo vale para las formas institucionales. La Caldera es una Comisión Municipal con
presidente designado por el gobernador en 1908, lo sigue siendo en 1922 —cuando el
Interventor Nacional nombra a sus cinco miembros y el Poder Ejecutivo le aprueba el
presupuesto— y todavía en 1983, en vísperas del traspaso democrático. En 2022 sigue
clasificada en la tercera categoría, como Vaqueros, que ya tiene población de segunda: La Caldera conserva la categoría que un decreto le fijó en 1931 y Vaqueros la que su ley de creación le dio en 1970, en un departamento de doce mil trescientos habitantes.

Y lo mismo vale para el problema del río. En enero de 1910 los vecinos del pueblo pidieron
al gobierno que se los defendiera «de las invasiones del Río del mismo nombre», y la
Provincia asignó fondos, nombró una comisión y le pidió al municipio que cofinanciara. En
febrero de 2026, diez días antes de que una jueza aprobara un plan de manejo, el municipio
informaba que construía bateas y muros de contención con maquinaria pesada.

De ahí la tercera tesis:

\begin{quote}
El dispositivo no se construyó para el mercado inmobiliario. El mercado inmobiliario
encontró un dispositivo ya armado —unidades, nombres, formas institucionales, obras de
defensa y opacidad dominial— y lo usó. Lo que cambió en el último cuarto de siglo no es el
territorio ni su administración: es el uso al que se destina el suelo.
\end{quote}

Esta tesis no reemplaza a las otras dos: las fecha. La regulación inversa y la opacidad del
criterio no son innovaciones de los años dos mil. Son formas antiguas de administrar un
territorio, que un ciclo inmobiliario reciente puso a trabajar en otra escala.

\subsection*{Lo que la tercera tesis no dice}

Formulada así, esa tesis describe una continuidad, y hay un punto en que el archivo la
contradice. Conviene decirlo aquí y no esconderlo en un capítulo interior, porque es la
objeción más seria que este libro se hace a sí mismo.

\textbf{En materia de agua, el municipio tuvo más poder en 1939 que en 2015.}

En 1914 el municipio de La Caldera dictó una ordenanza que reglamentaba la distribución de
\textbf{la totalidad} del caudal del río Santa Rufina, reconociendo derechos adquiridos por
uso tradicional. En 1939, al resolver el pedido que un comerciante de la capital había presentado en 1937 para obtener concesión sobre ese río
para regar la finca Campo Alegre, la Provincia \textbf{denegó} el pedido: hubo disenso
municipal, tres oposiciones de particulares, y el Poder Ejecutivo reconoció que otorgarla
alteraría el reparto fijado por aquella ordenanza. Al año siguiente confirmó la negativa
diciendo que \textbf{razones de orden legal le impedían} conceder agua de ese río.

Compárese con lo que el capítulo \ref{cap:ribera} documenta. Entre 2013 y 2015 la Secretaría de Recursos
Hídricos dictó veinticinco actos de línea de ribera sobre la cuenca del departamento, y en ninguno consta intervención municipal; la excepción, que la serie cuenta aparte, es la Resolución 023/14, dictada en 2014 a pedido del intendente sobre la ribera del propio pueblo. El capítulo \ref{cap:loteo} documenta que el plano de un loteo se
aprobó por resolución ministerial sin que figure ninguno de los diez actos que el
propio Código del municipio exige (apéndice \ref{cap:matriz}).

\textbf{Eso no es continuidad: es pérdida.} Y obliga a precisar la tesis en lugar de
abandonarla.

Lo que se mantiene idéntico durante un siglo son las \textbf{unidades} --- las fincas, los
nombres, los límites --- y la \textbf{forma institucional} del municipio: comisión de tercera
categoría, presidente designado, presupuesto aprobado desde afuera, renta que no alcanza. Lo
que cambió, y cambió en contra, es la \textbf{capacidad de decidir sobre los recursos de su
propio territorio}.

La tercera tesis queda entonces así:

\begin{quote}
El mercado inmobiliario encontró armado un dispositivo de unidades, nombres y formas
institucionales anterior a él. Pero no lo encontró intacto: encontró un municipio que había
perdido, en el camino, la voz que alguna vez tuvo sobre el agua y el suelo de su
jurisdicción.
\end{quote}

Cuándo y cómo se produjo esa pérdida es la pregunta que este libro no puede contestar del
todo, pero de la que tiene el punto de partida y una pista sobre el mecanismo. El punto de
partida documentado para La Caldera está en 1885: el capítulo \ref{cap:redes} documenta que en junio de ese año la Municipalidad de
La Caldera \textbf{concedió} treinta litros por segundo del río Wierna, y que el Estado
provincial reconoció esa facultad como legítima. El del mecanismo es anterior: en 1866 una ordenanza municipal de Campo Santo repartía el agua del Mojotoro y la Legislatura la sancionó como ley al año siguiente. La pista está en 1942, cuando ese mismo
derecho se \textbf{inscribe en los registros provinciales} y deja de constar en un acuerdo
municipal. \textbf{Y hay ahora un punto intermedio, que el barrido de 1918 aportó al cierre}: en
septiembre y en noviembre de ese año, dos decretos del Interventor Nacional nombran los Jueces de
Agua de los distritos de Caldera y de La Calderilla «\textit{vista la nota que antecede del Sr.\
Comisionado Municipal de Caldera, proponiendo \ldots\ las personas}» (capítulo
\ref{cap:siglo}). \textbf{La facultad municipal ya no es conceder: es proponer, y la firma es
provincial.} Conceder en 1885, proponer en 1918, inscribir en la provincia en 1942: la pérdida
tiene tres tiempos y el del medio recién aparece. El resto ocurre en los sesenta y dos años que el relevamiento no leyó de corrido ---y 1948 agrega un quinto tiempo, el más bajo: colaborar a pedido del organismo provincial (capítulo \ref{cap:siglo})---. El
capítulo \ref{cap:metodo} declara ese hueco entre los límites del trabajo, y el apéndice \ref{cap:pedidos} indica dónde
empezar a llenarlo.

\subsection*{Cómo se relacionan}

Las tres tesis describen el mismo objeto desde ángulos distintos, y conviene no confundirlas. Las dos primeras son sincrónicas: dicen cómo funciona el dispositivo. La tercera es de duración: dice desde cuándo.

La primera es sobre \textbf{a quién} se le exige. La segunda es sobre \textbf{qué se puede saber}. Y se articulan en un punto preciso: el solicitante de un lote fiscal tiene un baremo que lo mide, pero no puede consultarlo; el comprador de un lote privado tiene un reglamento que debería protegerlo, pero no puede leerlo. La opacidad no protege a los dos por igual, porque quien tiene abogado, tiempo y capacidad de gestión consigue el anexo, y quien no los tiene, no.

De modo que la segunda tesis no es un tema aparte: es el mecanismo por el cual la primera se sostiene sin que nadie tenga que decidirla.

\subsection*{Qué refutaría cada tesis}

Una tesis que no dice qué la desmentiría no se puede discutir: sólo se puede creer o no. Las tres se exponen aquí a la misma prueba que el libro aplica a cada expediente, con el documento que la resolvería y el lugar del apéndice \ref{cap:pedidos} donde ya está pedido.

\textbf{La primera tesis se cae si el circuito privado cumplió, en expediente, exigencias equivalentes a las del fiscal.} Su evidencia primaria es la Resolución 646D/21 frente a las Resoluciones 7, 8 y 9 de 2016 y al Anexo de la 82/16 (capítulos \ref{cap:loteo} y \ref{cap:tierrafiscal}). Si las fojas 58, 59, 68, 84, 97/101 y 137 del expediente 18-37.348/2019-0 contienen los certificados de no inundabilidad y de aptitud ambiental, la provisión efectiva de servicios y el visado hídrico que el anexo del Decreto 1682/19 exige, y si el registro de la Unidad Coordinadora muestra lo mismo para los demás loteos del corredor, la regulación no fue inversa: fue igual de exigente y este libro no la vio. Se cae también, por el otro extremo, si las actas de los sorteos y el registro público de tenedores precarios del artículo 1º inciso g de la Resolución 82/16 muestran que las exigencias escritas para el solicitante de tierra fiscal no se aplican en la práctica. Las dos series están pedidas a Inmuebles, al Boletín Oficial y a la Secretaría de Tierras y Bienes del Estado.

\textbf{La segunda tesis se cae si el criterio sí está escrito en algún lugar accesible.} Su evidencia primaria son los seis instrumentos del capítulo \ref{cap:opacidad} y los dos edictos de 2014 sobre la misma matrícula, uno con ciento trece pares de coordenadas y otro sin ninguno (capítulo \ref{cap:ribera}). Si la Secretaría de Recursos Hídricos tiene una resolución o instructivo publicado que fije cuándo se publican las coordenadas y cuándo no, o si los cinco anexos que no se obtuvieron se entregan a simple pedido, sin arancel y en plazo, y resultan no contener nada que decida ---ni obras exigidas, ni exenciones, ni puntos de ribera---, la opacidad es un defecto de infraestructura documental y no una forma de reservar el criterio. Los pedidos están en las secciones de Recursos Hídricos, del Boletín Oficial y de obra pública.

\textbf{La tercera tesis se cae si las unidades de hoy no descienden de las de 1909, o si el municipio recuperó la voz que perdió.} Su evidencia primaria son los deslindes de 1921 y 1932, el padrón de parajes de 1940, el revalúo de 1944 y la serie de 1885 a 1942 sobre el agua (capítulos \ref{cap:fincas}, \ref{cap:siglo} y \ref{cap:redes}). Si el antecedente dominial de las matrículas 4366 y 6.558 ---y de las demás matrículas de los loteos--- muestra que proceden de reagrupamientos posteriores y no de las fincas históricas, la continuidad de las unidades es sólo de nombres. Y si entre 1947 y 2012 ---el tramo que el relevamiento no leyó sobre la imagen: 1947 a 1950 sólo se barrieron, y de 1951 en adelante hubo consultas puntuales--- aparece un acto que devuelva al municipio competencia efectiva sobre el agua o el suelo de su jurisdicción, la «pérdida» que la tesis describe no es una pérdida sino una oscilación. Los folios y la serie de actos están pedidos a Inmuebles y al archivo provincial.

Ninguna de las tres condiciones se cumplió al cierre de este manuscrito. Ninguna fue puesta a prueba del todo: los documentos que las resolverían están pedidos y siguen sin obtenerse.

\section{Por qué este departamento y no otro}

La lámina \ref{fig:ubicacion} sitúa el departamento: el valle del río La Caldera baja de norte a
sur entre el dique Campo Alegre y Vaqueros, con la ruta nacional 9 como eje, y Vaqueros toca la
planta urbana del Gran Salta. Al oeste, las cumbres pasan de los cinco mil metros.

\begin{figure}[p]
\centering
\includegraphics[width=\textwidth,height=0.84\textheight,keepaspectratio]{fig-ubicacion.png}
\caption[El departamento La Caldera]{\textbf{El departamento La Caldera.} A: relieve, ríos, caminos, localidades y cerros del departamento, con el límite entre los municipios de La Caldera y Vaqueros y los parajes que registra la Base de Asentamientos Humanos (BAHRA). B y C: ubicación de la provincia en la Argentina y del departamento en la provincia. Elaboración propia con el Modelo Digital de Elevaciones MDE-Ar v2.1 de 30 metros y las capas SIG del Instituto Geográfico Nacional (hidrografía, red vial nacional, provincial y terciaria, huellas, municipios, plantas urbanas y cerros), la BAHRA, la cartografía de radios del Censo 2022 (INDEC) y los límites de la plantilla oficial del IGN para la Argentina parte continental americana. Elaboración propia; dibujo de \texttt{scripts/mapa\_ubicacion.py} del repositorio \texttt{mapas\_caldera} (apéndice \ref{cap:fuentes}).}\label{fig:ubicacion}
\end{figure}

\textbf{Primero: es la frontera, no el destino.} San Lorenzo es donde el patriciado salteño ya vive; La Caldera y Vaqueros son donde el mercado está \textit{produciendo} suelo urbano ahora. El proceso completo --- desmonte, apertura de camino, mensura, aprobación, escritura, preventa --- es observable en tiempo presente y con litigio activo.

\textbf{Segundo: es la cabecera hídrica de la capital, y lo es por expropiación.} En junio de 1972 un gobierno de facto declaró de utilidad pública seiscientas ochenta y seis hectáreas del departamento para construir el embalse Campo Alegre, con compensación fijada en valuación fiscal más treinta por ciento y toma de posesión previa a la discusión del precio. El acueducto que hoy lleva esa agua a la cisterna El Huaico recorre unos veinte kilómetros atravesando ambos municipios. La microcuenca cuyo bosque se remueve para lotear es la zona de recarga del sistema que abastece a la ciudad. El capítulo \ref{cap:expropiacion} reconstruye esa expropiación y la propone como punto de partida de todo lo demás.

\textbf{Tercero: hay expediente.} Existe una causa judicial --- CUIJ 17-00023860-0 --- que reúne en una sola pieza al municipio que autoriza, la provincia que omite, el desarrollador que ejecuta y los vecinos que pierden. Ninguna reconstrucción sociológica externa habría producido un material comparable al que produjeron las propias partes al defenderse.

\textbf{Cuarto: es la escala del fenómeno.} Entre 2010 y 2022 el departamento creció un 58,4 por ciento y fue el cuarto de mayor crecimiento demográfico de la Argentina, detrás de dos partidos del conurbano bonaerense y de Añelo, la capital de Vaca Muerta ---ordenamiento verificado sobre los resultados definitivos de veintitrés de las veinticuatro jurisdicciones, con la reserva que precisa el capítulo \ref{cap:poblacion}---. Lo administran dos municipios clasificados en tercera categoría.

\section{Las hipótesis heredadas y su saldo}

Del apéndice E de \textit{El dispositivo salteño}, retomadas aquí como programa:

\begin{description}[leftmargin=2.2cm,style=nextline]
\item[H.20] El mercado inmobiliario opera como mecanismo simultáneo de blanqueo de rentas no declaradas, consolidación intergeneracional de patrimonio y segregación residencial.
\item[H.21] Existe correlación entre ciclos de renta extractiva extraordinaria y desarrollo de emprendimientos de alta gama.
\item[H.22] Los apellidos que desarrollan y titularizan los emprendimientos se solapan con los documentados en el cuerpo principal del libro anterior.
\item[H.23] La cooperación entre escribanías, inmobiliarias y estudios jurídicos permite estructurar operaciones que no tributan como inversión sino como adquisición patrimonial.
\end{description}

Al cierre del relevamiento, el balance de estas cuatro hipótesis figura en el capítulo \ref{cap:conclusion}. Se adelanta lo principal: \textbf{la hipótesis del blanqueo, que motivó el proyecto, no encontró un solo elemento de prueba.}

\begin{cautela}
Ese saldo es magro y conviene declararlo de entrada. Un libro que se propuso verificar cuatro hipótesis y verifica, a lo sumo, dos en parte podría presentarse como fracaso. Lo que ocurrió, en cambio, es que el material desplazó el objeto dos veces: la pregunta por el blanqueo no encontró evidencia; la pregunta por la asimetría regulatoria, que nadie había formulado, encontró toda la que necesitaba; y la pregunta por la opacidad apareció sola, por repetición, en el propio trabajo de conseguir los documentos.

Se deja constancia porque es el modo honesto de exponer una investigación que cambió de tesis mientras se hacía.
\end{cautela}
